% ===================================================================
%  GAMBAR No. 9 — Gunung lan Kutha Panas. --- Alas lan Mbebedhag.
%
%  Traduction, faite en 2026, en javanais d'aujourd'hui, sur l'ido de
%  texto/io. Regles de la colonne : voir 10-gambar-01.tex. Deux
%  scenes, deux numerotations. Un bloc marque « ¶2 » par ordo.py.
%
%  LA MONTAGNE EST LE TABLEAU OU LE JAVANAIS EST LE PLUS A L'AISE DU
%  LIVRET, ET C'EST INATTENDU. Java est une ile de volcans, et le
%  vocabulaire y est complet et propre : gunung la montagne, punthuk
%  la colline, jurang le ravin (32), kawah le cratere, wedhus gembel
%  la nuee ardente, lahar la coulee, grojogan la cascade (10), tuk la
%  source, sumber, kali, lebak la vallee (25), tebing la paroi,
%  gua la grotte (56), watu, wedhi. Le tableau 9 dit un paysage des
%  Pyrenees, et la colonne l'ecrit sans emprunter un mot pour la
%  moitie des objets.
%
%  ET « lahar » (la coulee de lave) EST SORTI DE JAVA, COMME wajan AU
%  TABLEAU 6. Le mot est javanais, il a passe au malais, puis a
%  l'anglais et au francais des volcanologues : un lahar se dit lahar
%  a Paris comme a Yogyakarta. C'est le quatrieme mot du releve qui
%  voyage dans ce sens-la — قلی, بازار, wajan, lahar —, et le seul des
%  quatre qui soit parti pour une raison scientifique et non
%  commerciale.
%
%  MAIS LA NEIGE ET LA GLACE MANQUENT ENCORE, ET DEUX FOIS. Les
%  glaciers (5) et les sommets neigeux (4) demandent « gletser » et
%  « salju », comme au tableau 7 ; l'avalanche n'a aucun mot et
%  s'ecrit « longsoran salju », glissement de neige — sur le modele de
%  « longsor », le glissement de terrain, que Java connait trop bien.
%  Une langue qui n'a pas la chose emprunte le nom ; une langue qui a
%  une chose VOISINE compose. Les deux procedes sont dans le meme
%  alinea.
%
%  L'OURS (44) DU BOHEMIEN ET LE CERF (50) DE LA CHASSE A COURRE SONT
%  DEUX ANIMAUX QUE JAVA N'A PAS, ET ILS SE COMPORTENT AUTREMENT L'UN
%  QUE L'AUTRE. « bruwang » l'ours est un mot malais ancien, parce
%  qu'il y a des ours a Sumatra et a Borneo, tout pres ; « rusa » le
%  cerf est malais aussi, et le cerf de Java existe. Le sanglier (21)
%  se dit « celeng », mot javanais pur, et l'ile en est pleine. C'est
%  donc la faune d'EUROPE, ici, qui est locale par accident : chamois
%  et marmotte manquent, le reste est chez lui.
%
%  LE CYGNE (45) ET L'OIE DU TABLEAU 7 SE PARTAGENT LE MEME OISEAU EN
%  DEUX MOTS DE DEUX ORIGINES. L'oie se dit « banyak », mot javanais
%  qu'on avait du oter de kolonoj.py au tableau 3 ; le cygne se dit
%  « angsa », du sanskrit haṃsa, le hamsa des temples. L'oiseau qu'on
%  eleve dans la cour a garde son nom d'ici, l'oiseau qu'on ne voit
%  que sur une frise a pris celui de l'Inde. Une premiere version
%  ecrivait « banyak » pour le cygne : la faute est corrigee, et elle
%  etait instructive — ce sont les deux mots les plus proches du
%  tableau et ils viennent d'aussi loin l'un de l'autre que possible.
%
%  ET LE PIC-VERT (30) EST LE SEUL OISEAU DU TABLEAU QUI RESISTE.
%  Java a des pics, mais aucun nom courant ne les separe des autres
%  grimpeurs ; la colonne ecrit « manuk pelatuk », qui est le nom
%  javanais du pic et qui veut dire « l'oiseau qui frappe ». Il decrit
%  le geste, comme pandeleng et pangrungu decrivent les sens au
%  tableau 2 : c'est la maniere ordinaire de cette langue de nommer.
% ===================================================================

%%K t09-apar-1 apar
\VUsaut{4.0mm}
\VUtitre{15.0pt}{60}{GAMBAR No. 9}
\VUsaut{3.0mm}

%%K t09-tit-1 sub
\VUsaut{3.0mm}
\VUcentre{12.0pt}{0}{\textit{Adegan Kapisan.}}
\VUsaut{3.0mm}

%%K t09-tit-2 sub
\VUsaut{3.0mm}
\VUpk{11.4pt}{40}{Gunung. --- Kutha Panas.}
\VUsaut{3.0mm}

%%K t09-c1-01-1 p
1. --- Wis wolung dina aku ana ing Pratz. Aku ora bisa ngucapake kabeh
kaguman kang dakrasakake nalika aku ngadeg ing ngarepe
\VUgras{pagunungan} \textsuperscript{(1)} Pirenei kang nggumunake, kang
\VUgras{gigire} \textsuperscript{(2)} katon ing langit biru kaya rerenggan
raseksa.

%%K t09-c1-02-1 p
2. --- Aku ora ngira yen \VUgras{pucuk} \textsuperscript{(3)} Pirenei
tekan \VUgras{dhuwur} kang semono gedhene, lan aku gumun banget ing
ngarepe \VUgras{gletser} \textsuperscript{(5)} kang endah banget lan
\VUgras{pucuking gunung} \textsuperscript{(4)} kang salju-salju, panggonan
\VUgras{longsoran salju} kang medeni temen padha nggelundhung.

%%K t09-tit-3 sub
\VUsaut{3.0mm}
\VUpk{10.2pt}{40}{Pemandangan gunung.}
\VUsaut{3.0mm}

%%K t09-c1-03-1 p
3. --- Sanalika ing dina sawise tekaku, aku menyang \VUgras{gunung
geni} \textsuperscript{(6)} lawas kang wis mati kang ana ing cedhake
Pratz. Ing pinggire \VUgras{kawah} kang garing lan gundhul, wong isih
weruh cetha tilase liwate \VUgras{lahar} kang mili, pirang-pirang abad
kepungkur, ing \VUgras{ereng-erenge gunung} \textsuperscript{(7)}. Aku
bisa nemu, ing salah sijine \VUgras{ereng-erengan,} sawetara
gempilaning lahar iku.

%%K t09-c1-04-1 p
4. --- Pratz manggon ing tapel wates, lan \VUgras{beteng}
\textsuperscript{(8)} kang njaga \VUgras{lengkehan} \textsuperscript{(27)}
kang misahake Prancis karo Spanyol pancen ngelingake marang kastil
lawas ing pinggire Rhein kang katon ngintip \VUgras{buron} saka pucuke
watu.

%%K t09-c1-05-1 p
5. --- \VUgras{Manuk garudha} \textsuperscript{(9)} kang gedhe banget,
kang duwe susuh ora adoh saka \VUgras{beteng} \textsuperscript{(8)},
kerep narik kawigatenku.

Manuk iku \VUgras{manuk pemangsa} kang \VUgras{cucuk}e kuwat, kang
kerep nggawa cempe utawa anak wedhus jawa marang anake, kang dicekel
nganggo \VUgras{kuku}ne. Semono gedhene \VUgras{elar}e nganti bisa
nglayang suwe karo momotan kang abot iku.

%%K t09-tit-4 sub
\VUsaut{3.0mm}
\VUpk{10.2pt}{40}{Wong Plesir.}
\VUsaut{3.0mm}

%%K t09-c1-06-1 p
6. --- Ing kono mlaku-mlaku kang paling nyenengake yaiku plesir menyang
grojogan \textsuperscript{(10)} « Kau ». \VUgras{Grojogan} iku tiba
mubeng-mubeng banjur ilang dadi \VUgras{kali cilik} kang endah ing
\VUgras{alas cemara} \textsuperscript{(11)} ing sisih kulone kutha. Kita
munggah liwat alas iku tekan \VUgras{papan pengamatan cuaca}
\textsuperscript{(12)}, lan saka pucuking \VUgras{panggonan nyawang,}
kita weruh \VUgras{panorama} tlatah iku kabeh. \VUgras{Pemandangane}
nggumunake, lan \VUgras{alam} kaya-kaya ngesokake marang tanah iku
kabeh raja brana kang diduweni.

%%K t09-c1-07-1 p
7. --- Kanggo mudhun, kita numpak \VUgras{kreta kabel}
\textsuperscript{(13)}, lan kita mandheg ing \VUgras{setasiun} cilik
cedhak \VUgras{jugrugan} \textsuperscript{(14)} \VUgras{kastil feodal}
lawas kang biyen dienggoni para bendarane Pratz.

%%K t09-c1-08-1 p
8. --- Mesakake temen kastil iku! Apa kira-kira kang dipikir nalika
nyawang ing ngisor \VUgras{kutha panas} \textsuperscript{(15)} kang apik,
kang saiki dadi \VUgras{papan pengobatan panas} kang lagi rame?

%%K t09-c1-08-2 p
\VUgras{Hawa}ne kepenak temen, \VUgras{banyu mineral}e mandi temen,
nganti \VUgras{pengobatan} kang dilakoni ing \VUgras{sanatorium}
\textsuperscript{(17)} dadi kesenengan kang tenanan.

%%K t09-tit-5 sub
\VUsaut{3.0mm}
\VUpk{10.2pt}{40}{Kutha panas kang nyenengake.}
\VUsaut{3.0mm}

%%K t09-c1-09-1 p
9. --- Kejaba iku, wong wis nindakake sakabehe supaya manggon ing kono
kepenak. \VUgras{Kreteg dhuwur} \textsuperscript{(16)} kang gedhe wis
dibangun supaya sepur bisa liwat ; \VUgras{taman}
\textsuperscript{(18)} \VUgras{papan pengobatan panas,} panggonan wong
\VUgras{nganakake konser,} ditata kanthi endah. Sawetara
« \VUgras{vila} » \textsuperscript{(19)} kang apik wis dibangun ing
sakubenge kasino \textsuperscript{(21)}, lan \VUgras{dalan gedhe}
\textsuperscript{(22)} kang dirumat becik banget nggampangake
sesambungan.

%%K t09-c1-10-1 p
10. --- \VUgras{Tanggul} \textsuperscript{(23)} kang prasaja misahake
\VUgras{dalan} \textsuperscript{(20)} karo \VUgras{lebak}
\textsuperscript{(25)} kang sabenere, kang ing dhasare ana
\VUgras{tlaga} \textsuperscript{(26)} cilik kang endah kang mbanyoni
\VUgras{kali cilik} \textsuperscript{(24)} kang bening.

%%K t09-c1-11-1 p
11. --- Ing sisih sijine lebak, \VUgras{tambang wesi}
\textsuperscript{(28)} lagi digarap. Kaping bola-bali aku nyawang
\VUgras{tukang tambang} mudhun ing \VUgras{sumuran,} lan aku malah
mlebu \VUgras{lorong} panggonan wong-wong iku nglampahi dinane, karo
njupuki \VUgras{bijih} kang isi \VUgras{logam.}

%%K t09-c1-12-1 p
12. --- \VUgras{Pabrik peleburan} kang gedhe wis dibangun cedhak
tambang supaya raja brana kang dijupuk saka kono sanalika bisa
digunakake. \VUgras{Tungku gedhe} \textsuperscript{(29)}, kang
dipanasi \VUgras{areng watu} kang akeh ing kene, mbisakake wong enggal
lan murah oleh \VUgras{wesi tuwang.}

%%K t09-c1-13-1 p
13. --- Ora adoh saka tambang, wong ndhudhuk \VUgras{tambang watu}
\textsuperscript{(30)}. Dudu \VUgras{granit,} dudu \VUgras{porfir,}
dudu uga \VUgras{watu pasir} kang dipara \VUgras{tukang tambang watu}
tuwa nganggo \VUgras{linggis}e, nanging mung \VUgras{watu tembok}
kanggo mbangun omah.

%%K t09-c1-14-1 p
14. --- Salah sijine rerenggan kang paling endah ing tanah kono yaiku
\VUgras{wit cemara} \textsuperscript{(31)}. Ing ngendi-endi ing dhasare
\VUgras{jurang} \textsuperscript{(32)}, ing pinggire \VUgras{grojogan
deres} \textsuperscript{(33)}, \VUgras{jurang jero}
\textsuperscript{(34)} utawa pereng, dom-dome kang lembut menehi warna
ijo.

%%K t09-tit-6 sub
\VUsaut{3.0mm}
\VUpk{10.2pt}{40}{Mlaku sikil.}
\VUsaut{3.0mm}

%%K t09-c1-15-1 p
15. --- Aku nganggo preinanku kanggo \VUgras{mlaku-mlaku adoh.}
\VUgras{Dalan} \textsuperscript{(35)} kang mendhat-mendhut ing gunung
mbisakake nempuh, tanpa krasa adohe, pirang-pirang \VUgras{kilometer}
lan kerep pirang-pirang \VUgras{puluhan kilometer.}

%%K t09-c1-15-2 p
\VUgras{Watu wates} \textsuperscript{(36)} lan \VUgras{cagak pratandha}
\textsuperscript{(37)} menehi tenger ing dalan, panggonan kita kerep
kepethuk \VUgras{wong gunung} \textsuperscript{(38)} enom karo
\VUgras{tas rangkep} \textsuperscript{(40)} ing iringane, kang nggawa
\VUgras{marmot} \textsuperscript{(39)}, utawa sawetara \VUgras{wong
gipsi} \textsuperscript{(41)} kang \VUgras{sarung guluné wulu}
\textsuperscript{(42)} aneh temen bedane karo \VUgras{kupluk}e
\textsuperscript{(43)} lawas kang wis suwek. Wong iku nuntun
\VUgras{bruwang} \textsuperscript{(44)} kang wis dilatih nganggo
\VUgras{ranté.}

%%K t09-tit-7 sub
\VUsaut{3.0mm}
\VUpk{10.2pt}{40}{Munggah gunung.}
\VUsaut{3.0mm}

%%K t09-c1-16-1 p
16. --- Meh saben dina aku lunga karo \VUgras{pemandu}
\textsuperscript{(48)} nglakoni \VUgras{munggah,} lan aku bangga banget
yen, karo \VUgras{tas} \textsuperscript{(47)} ing geger lan
« \VUgras{alpenstock} » ing tangan, kaya \VUgras{wong plesir}
\textsuperscript{(46)} kang tenanan, aku nyerbu \VUgras{gugusan watu}
\textsuperscript{(45)}.

%%K t09-c1-17-1 p
17. --- Saka pucuking gunung, wong-wong ing dhataran katon ora luwih
gedhe tinimbang semut ; \VUgras{pepanthan wedhus}
\textsuperscript{(50)} kang nyenggut ing \VUgras{pangonan}
\textsuperscript{(49)} katon kaya dolanane bocah. \VUgras{Wit pinus}
\textsuperscript{(51)} utawa uga \VUgras{omah kayu} \textsuperscript{(52)}
katon mung kaya cletheke ing ijo-ijoan.

%%K t09-c1-18-1 p
18. --- Sajrone plesiran iku, kita kerep kepethuk ing \VUgras{dalan
cilik} \textsuperscript{(53)} karo \VUgras{tukang mbebedhag kambing
gunung} \textsuperscript{(54)} kang manggul kewan kang lagi wae
dipateni.

%%K t09-c1-19-1 p
19. --- Aku nyelehake ing herbariumku pange \VUgras{pakis}
\textsuperscript{(55)} kang endah temen lan pang cilik \VUgras{eres}
\textsuperscript{(57)} kang abang mawar, kang dakpethik ing lawange
\VUgras{gua} \textsuperscript{(56)} cilik nalika mlaku-mlakuku kang
pungkasan.

%%K t09-tit-8 sub
\VUsaut{3.0mm}
\VUcentre{12.0pt}{0}{\textit{Adegan Kapindho.}}
\VUsaut{3.0mm}

%%K t09-tit-9 sub
\VUsaut{3.0mm}
\VUpk{11.4pt}{40}{Alas. --- Mbebedhag.}
\VUsaut{3.0mm}

%%K t09-c2-01-1 p
1. --- Wingi aku nglampahi sedina kang nyenengake banget ing alas.
Kesregepan kang katon ing antarane kewan lan \VUgras{serangga}
\textsuperscript{(5)} kang manggon ing kono narik kawigatenku banget.

%%K t09-c2-01-2 p
\VUgras{Kemangga} \textsuperscript{(1)} kang nganam \VUgras{jaring}e
\textsuperscript{(2)} ing antarane pang loro kanggo nyekel \VUgras{laler}
\textsuperscript{(3)} kang liwat, \VUgras{keong} \textsuperscript{(4)}
kang nggawa omahe kanthi rekasa, kumbang tandhuk kang nggoleki
pangane, semut kang sregep lan mempeng ing cedhake \VUgras{omah semut}
\textsuperscript{(6)} — kabeh kewan cilik iku lunga, teka, nyambut gawe
kanthi sregep kang ora lumrah.

%%K t09-c2-02-1 p
2. --- \VUgras{Ula} \textsuperscript{(7)} kang sepisanan dakkira
\VUgras{ula bandhotan,} nanging kang jebul mung \VUgras{ula sawah}
lumrah, ndhelik ing sangisore watange wit, lan sedhela-sedhela
ngetokake sirahe kang lancip kang tansah katon siyap nyokot. Ing
ngarepku \VUgras{silit kebo} \textsuperscript{(8)} kang gedhe ngrayap
kanthi abot sawise mangan salah sijine \VUgras{arben}
\textsuperscript{(11)} kang enak kang thukul akeh ing kono.

%%K t09-c2-03-1 p
3. --- Lemahe kegelaran \VUgras{kembang alas} : \VUgras{kembang primula}
\textsuperscript{(9)}, \VUgras{kembang pervinka} \textsuperscript{(10)},
lan akeh tanduran liya kang ora dakkenal, padha kembang ing antarane
\VUgras{gerumbul suket} \textsuperscript{(12)}.

%%K t09-c2-04-1 p
4. --- Ing tlatah kono, \VUgras{truffle} meh padha lumrahe karo
\VUgras{jamur} \textsuperscript{(14)}, lan \VUgras{gembluk}
\textsuperscript{(13)} duweke tukang jaga alas kanthi rakus nggoleki
apa dheweke bisa nemu sawetara.

%%K t09-tit-10 sub
\VUsaut{3.0mm}
\VUpk{10.2pt}{40}{Mbebedhag colongan.}
\VUsaut{3.0mm}

%%K t09-c2-05-1 p
5. --- Aku nyekseni \VUgras{pungkasane} adegan kang nyenengake banget.
Kanthi bantuan pangambune \VUgras{asu basset} \textsuperscript{(15)},
\VUgras{tukang jaga alas} \textsuperscript{(16)} lagi wae ngagetake
\VUgras{tukang mbebedhag colongan} \textsuperscript{(22)}, pas nalika
wong iku tata-tata nggawa \VUgras{buron} \textsuperscript{(17)} kang wis
dipateni. Luwih milih ninggal buruane tinimbang kecekel, tukang
mbebedhag colongan iku wis mlayu, lan \VUgras{terwelu alas}
\textsuperscript{(18)}, \VUgras{rusa wadon} \textsuperscript{(19)},
\VUgras{kidang} \textsuperscript{(20)} lan \VUgras{celeng}
\textsuperscript{(21)} gumlethak ing suket, gawe bungahe tukang jaga
alas.

%%K t09-c2-06-1 p
6. --- Maneka warnane wit kang ana ing alas iku nggumunake. \VUgras{Wit
beech} \textsuperscript{(23)} lan \VUgras{wit birch} \textsuperscript{(26)},
\VUgras{wit pinus,} kang ngasilake \VUgras{damar} \textsuperscript{(27)},
\VUgras{wit jati} \textsuperscript{(29)} lan uga akeh liyane nyampurake
pang-pangane.

%%K t09-c2-06-2 p
\VUgras{Oyot rambat} \textsuperscript{(24)}, kang nemplek ing watange wit
dhuwur, menehi warna ijo mencorong marang \VUgras{alas}
\textsuperscript{(25)}, kang nyawiji kanthi becik karo warna kang luwih
padhang lan ijo enom saka \VUgras{lumut} \textsuperscript{(31)}
panggonan \VUgras{manuk pelatuk} \textsuperscript{(30)} nggoleki
serangga.

%%K t09-tit-11 sub
\VUsaut{3.0mm}
\VUpk{10.2pt}{40}{Pemandangan alas.}
\VUsaut{3.0mm}

%%K t09-c2-07-1 p
7. --- Wit jati lawas gawe kesengsemku. \VUgras{Oyot}e
\textsuperscript{(32)} kang gedhe lan mulet, separo metu saka lemah,
menehi wujud kang endah temen bab kekuwatan lan kasantosan, lan aku
takon marang awakku dhewe kepriye \VUgras{sing duwe}
\textsuperscript{(35)} bisa lila negor lan masrahake wit iku marang
\VUgras{graji} \textsuperscript{(34)} \VUgras{tukang negor}
\textsuperscript{(33)}. \VUgras{Watang} \textsuperscript{(36)} kang
mesakake, kang wis dicopoti \VUgras{kulite} \textsuperscript{(37)}
utawa dibelah \VUgras{wadung} \textsuperscript{(38)} \VUgras{buruh
harian} \textsuperscript{(39)}, gawe sedhihku.

%%K t09-c2-08-1 p
8. --- Ing cedhake, \VUgras{tukang gawe areng} \textsuperscript{(41)}
tuwa wis nyiyapake \VUgras{tumpukan} \textsuperscript{(42)} areng nganggo
sekope, lan wong wadon tuwa nggawa \VUgras{bongkokan}
\textsuperscript{(40)} \VUgras{kayu garing} saka pang-pange wit kang wis
ditegor.

%%K t09-tit-12 sub
\VUsaut{3.0mm}
\VUpk{10.2pt}{40}{Mbebedhag.}
\VUsaut{3.0mm}

%%K t09-c2-09-1 p
9. --- Aku arep ngaso sedhela ing \VUgras{omahe tukang jaga alas}
\textsuperscript{(46)}, nanging, nalika aku tekan cedhake \VUgras{tlaga
cilik} \textsuperscript{(43)}, panggonan \VUgras{angsa} \textsuperscript{(45)}
kang endah nglangi, aku weruh yen \VUgras{banjir}
\textsuperscript{(44)} kang lagi wae kelakon wis ngowahi dalan iku dadi
\VUgras{rawa} tenan. Aku kudu minger, lan, nalika liwat cedhak
\VUgras{wit cedar} \textsuperscript{(49)}, \VUgras{wit poplar}
\textsuperscript{(48)} lan \VUgras{wit elm} \textsuperscript{(47)} kang
ana ing pinggire alas, aku weruh metu saka \VUgras{grumbulan}
\textsuperscript{(54)} \VUgras{rusa lanang} \textsuperscript{(50)} kang
endah banget, kang diburu \VUgras{pepanthan asu buron}
\textsuperscript{(51)} kang ora gelem lerem, kang uga disemangati
\VUgras{sungu} \textsuperscript{(53)} \VUgras{wong mbebedhag kang
nunggang jaran} \textsuperscript{(52)}. Kewan kang mesakake iku wis
kentekan tenaga. Rusa iku tiba mati ing sawetara jangkah saka aku.

%%K t09-c2-10-1 p
10. --- Anak lanange tukang jaga alas, kang katarik dening rame-ramene
\VUgras{mburu,} metu saka omah, lan kita wong loro bali menyang alas.
Dheweke wis weruh \VUgras{manuk murai} \textsuperscript{(56)}
ing pange wit jati lawas. Dheweke ngira yen susuhe mesthi ndhelik ing
\VUgras{godhongan} \textsuperscript{(58)}, lan dheweke kepengin njupuk
susuh iku.

%%K t09-c2-10-2 p
Cukat kaya \VUgras{bajing} \textsuperscript{(61)}, dheweke menek saka
\VUgras{pang} \textsuperscript{(55)} menyang pang, nanging ora nemu
apa-apa lan mung bisa gawe mabure \VUgras{manuk gagak}
\textsuperscript{(60)} tuwa kang nyanggong ing \VUgras{ranting}
\textsuperscript{(57)}.

\VUsaut{4.0mm}
\VUfilet{22mm}
